\section{Арифметика предела функции. Предельный переход в неравенствах. Теорема о зажатой функции. (Beta)}

\subsection{Арифметика предела функции}

\begin{theorem}
Пусть $f(x)$ и $g(x)$ заданы в проколотой окрестности точки $x_0$ и
\[
    \Lim{x}{x_0} f(x) = A,\qquad \Lim{x}{x_0} g(x) = B.
\]
Тогда выполняются следующие утверждения:
\begin{enumerate}
    \item $\Lim{x}{x_0} \bigl( f(x) + g(x) \bigr) = A + B$;
    \item $\Lim{x}{x_0} \bigl( f(x) - g(x) \bigr) = A - B$;
    \item $\Lim{x}{x_0} \bigl( f(x)\cdot g(x) \bigr) = A\cdot B$;
    \item Если $B\neq 0$ и $g(x)\neq 0$ в проколотой окрестности $x_0$, то
    \[
        \Lim{x}{x_0} \frac{f(x)}{g(x)} = \frac{A}{B}.
    \]
    \item Если для некоторого целого $k\ge 1$ имеется дополнительное требование (при необходимости: $f(x)\ge 0$ при $x$ близко к $x_0$ для нечётных корней и т.\,п.), то
    \[
        \Lim{x}{x_0} \sqrt[k]{f(x)} = \sqrt[k]{A}.
    \]
\end{enumerate}
\end{theorem}

\begin{Proof}
Доказательство удобно проводить через определение по Гейне (через последовательности).

Возьмём произвольную последовательность $\{x_n\}_{n\in\N}$ такая, что
$x_n\to x_0$ и $x_n\neq x_0$ для всех $n$. Тогда по условию
$f(x_n)\to A$ и $g(x_n)\to B$ при $n\to\infty$. По известным свойствам пределов последовательностей имеем:
\[
f(x_n)+g(x_n)\to A+B,\quad f(x_n)-g(x_n)\to A-B,
\]
\[
f(x_n)g(x_n)\to AB,\qquad
\text{и при }B\neq0:\ \frac{f(x_n)}{g(x_n)}\to\frac{A}{B}.
\]
Применяя обратное утверждение определения по Гейне, получаем соответствующие предельные переходы для функций при $x\to x_0$.

Остальные пункты (корень и т.\,п.) доказываются аналогично с учётом непрерывности элементарных отображений в точках соответствующих значений.
\end{Proof}

\begin{example}
Покажем на примере:
\[
\Lim{x}{0}\frac{\sin x + x^2}{x} = \Lim{x}{0}\frac{\sin x}{x} + \Lim{x}{0} x = 1 + 0 = 1.
\]
Здесь использована арифметика предела и известный предел $\Lim{x}{0}\frac{\sin x}{x}=1$.
\end{example}

\subsection{Предельный переход в неравенствах}

\begin{theorem}
Пусть в некоторой проколотой окрестности точки $x_0$ выполнено неравенство
\[
    f(x) \le g(x)\qquad \text{для всех }x\text{, достаточно близких к }x_0,
\]
и существуют пределы $A=\Lim{x}{x_0} f(x)$ и $B=\Lim{x}{x_0} g(x)$. Тогда
\[
    A \le B.
\]
\end{theorem}

\begin{Proof}
Опять воспользуемся определением по Гейне. Возьмём любую последовательность
$\{x_n\}$, $x_n\to x_0$, $x_n\neq x_0$. По условию для всех достаточно больших $n$
выполняется $f(x_n)\le g(x_n)$. При проходе к пределу получаем $A\le B$,
поскольку $f(x_n)\to A$ и $g(x_n)\to B$.
\end{Proof}

\begin{remark}
Аналогично: если $f(x)<g(x)$ в проколотой окрестности и оба предела существуют, то $A\le B$, причём строгое неравенство $A<B$ в общем случае может не выполняться (пример: $f(x)=x,\ g(x)=x+1/n$ по подходящим последовательностям).
\end{remark}

\subsection{Теорема о зажатой (сжатой) функции}

\begin{theorem}[Теорема о зажатой функции]
Пусть функции $f,g,h$ определены в проколотой окрестности точки $x_0$ и существует $\delta_0>0$ такое, что
\[
    f(x)\le h(x)\le g(x)\qquad \text{при }0<|x-x_0|<\delta_0.
\]
Если
\[
    \Lim{x}{x_0} f(x) = \Lim{x}{x_0} g(x) = A,
\]
то существует предел $\Lim{x}{x_0} h(x)$ и он равен $A$.
\end{theorem}

\begin{Proof}
Возьмём произвольную последовательность $\{x_n\}$ с $x_n\to x_0$, $x_n\neq x_0$.
Для всех достаточно больших $n$ имеем $f(x_n)\le h(x_n)\le g(x_n)$. Переходя к пределу и используя, что $f(x_n)\to A$ и $g(x_n)\to A$, получаем по теореме о зажатой для последовательностей:
\[
h(x_n)\to A.
\]
Так как это верно для любой такой последовательности, по определению по Гейне $\Lim{x}{x_0} h(x)=A$.
\end{Proof}

\begin{consequence}
Если $f(x)\le h(x)\le g(x)$ и $\Lim{x}{x_0} f(x)=\Lim{x}{x_0} g(x)=0$, то $h(x)=\overline{o}(1)$ при $x\to x_0$.
\end{consequence}

\begin{example}
Докажем классический пример:
\[
    \Lim{x}{0} x^2\sin\frac{1}{x} = 0.
\]
Заметим, что $-1\le \sin\frac{1}{x}\le 1$, откуда
\[
    -x^2 \le x^2\sin\frac{1}{x} \le x^2.
\]
Пределы боковых функций при $x\to0$ равны нулю, значит по теореме о зажатой предел искомой функции равен $0$.
\end{example}

\begin{task}
Найти предел
\[
    \Lim{x}{0} \frac{x-\sin x}{x^3}.
\]
\end{task}

\begin{solution}
Разложим $\sin x$ в ряд Тейлора:
\[
\sin x = x - \frac{x^3}{6} + o(x^3),\quad x\to0.
\]
Отсюда
\[
\frac{x-\sin x}{x^3} = \frac{x - \bigl(x - \frac{x^3}{6} + o(x^3)\bigr)}{x^3}
= \frac{1}{6} + \frac{o(x^3)}{x^3} \to \frac{1}{6}.
\]
Следовательно предел равен $\dfrac{1}{6}$.
\end{solution}

\begin{hint}
Если вы ещё не знакомы с рядом Тейлора, можно решить задачу через правило Лопиталя (дважды применив) или через известные разложения для $\sin x$.
\end{hint}
